\chapter{QAOA}
Nei capitolo precedenti ci siamo concentrati sull'ottenere una Hamiltoniana il cui ground state rappresentasse la soluzione del nostro problema. Implicitamente o espicitamente facevamo riferimento ai quantum annealer come macchine per trovare la configurazione che minimizzasse l'energia.

In questo capitolo mostreremo come passare da un problema formulato per quantum annealer a un problema formulato per computer quantum gate based attraverso l'algoritmo \emph{Quantum Approximation Optrimization Algoritm} (QAOA) Proposto per la prima volta da Farhi et al. in~\cite{bibid}.

Iniziamo descrivendo l'algoritmo QAOA evidenziando vantaggi e svantaggi teorici rispetto a QAC\footnote{Quantum Adiabatic Computing (Sezione~\ref{})} . Mostreremo poi come sia possibile tradurre un problema formulato in forma QUBO (Sezione~\ref{}), ottenuto al termine dalla pipeline di QA-Prolog, in un problema adatto ai quantum gate based computer di IBM.

Infine testeremo la traduzione proposta e mostreremo alcuni risultati sperimentali discutendoli.
\section{QAOA}
L'algoritmo QAOA è ispirato al QAC, entrambi tentano di trovare l'uatovalore minimo di una matrice che rappresenta l'energia di un sistema fisico~\cite{fakhimi2023quantum}. Questa matrice descrive un osservabile, di conseguenza è un operatore Hermitiano e in particolare è l'Hamiltonana $\mathbf{H}$ del sistema. L'autovalore più piccolo di $\mathbf{H}$ è l'energia minima del sistema che si può essere misurata per questo è detta ground state energy.

Come descritto nella Sezione~\ref{} il QAC viene eseguito su un quantum annealer che parte da una Hamiltonana di cui e facile conoscere il graound state ed evolve, adiabaticamente (per rimanere nel ground state) verso l'Hamiltoniana con ground state ignoto. Il tempo di evoluzione da un'Hamiltonana all'altra può essere esponenziale\cite{farhi2014quantum}, inoltre generalmente la probabilità di successo di QAC non aumenta all'aumentare del tempo di annealing.

Il QAOA targetta macchine quantum gate based e una volta ottimizzati i suoi parametri la bontà del risultato crese all'aumentare del numero di iterazioni, dette \emph{livelli}.

\subsection{Circuito}
L'algoritmo QAOA comincia con uno stato $\Ket{s}$ con sovrapposizione massima. È possibile ottenere questo stato partendo dallo stato con $n$ \emph{qubit} a zero ($\Ket{0}$) e applicando ad ogniuno la porta di Hadamard (H)~\cite{fakhimi2023quantum}, denotiamo questo operatore con $H^{\otimes n}$. 

Il circuito è poi costituito da $p$ livelli tutti uguali. Ogni livello è composto da due operatori: 
\begin{itemize}
	\item phase separation operator: $U_P(\gamma) = e ^{-i\gamma\mathbf{H}}$ dove $\mathbf{H}$ è l'amiltoniana del sistema e $\gamma \in [0, 2\pi]$ è un parametro;
	\item mixing operator:  $U_M(\beta) = e ^{-i\beta B}$ dove B è l'operatore che applica il Pauli Gates X a ogni \emph{qubit} e $\beta \in [0, \pi]$ è un parametro.
\end{itemize}
L'operatore $U_(P\gamma)$ simula $\mathbf{H}$, mentre $U_M(\beta)$, che non commuta con $U_P(\gamma)$\footnote{$U_P(\gamma) \times U_M(\beta) \neq U_M(\beta) \times U_P(\gamma)$.}, permette di modificare la distribuzione di probalbilità dell'output~\cite{fakhimi2023quantum}. 

Usare i Pauli Gates X non è l'unico modo per costriure il mixing operator, in \cite{wang2020xy, hadfield2019quantum} si esplorano altri mixing operetor.

\begin{figure}[h]
	\centering
	\includegraphics[width=.8\linewidth]{qaoa/qaoa}
	\caption{\cite{blekos2024review}}
	\label{fig:qaoa}
\end{figure}

La Figura~\ref{fig:qaoa} mostra lo schema generale dell'algoritmo QAOA considerando $p$ livelli.

Chiamati $\vec{\gamma} = (\gamma_1, \gamma_2,\dots,\gamma_p)$ e$\vec{\beta} = (\beta_1, \beta_2,\dots,\beta_p)$, indichiamo con $\Ket{\psi(\vec{\gamma},\vec{\beta})}$ the \emph{ansatz} (una stima della soluzione) del QAOA. L'ansatz è definita come:
\[
	\Ket{\psi(\vec{\gamma},\vec{\beta})} = U_M(\beta_p)U_P(\gamma_p)\cdots U_M(\beta_1)U_P(\gamma_1)\Ket{s}
\]
Il circuito termina con la misurazione, il valore atteso di questa misura è:
\[
F(\vec{\gamma},\vec{\beta}) =\Braket{\psi(\vec{\gamma},\vec{\beta}) | \mathbf{H} | \psi(\vec{\gamma},\vec{\beta})}
\]

$F(\vec{\gamma},\vec{\beta})$ rappresenta il valore atteso di $\mathbf{H}$ nello stato $\Ket{\psi(\vec{\gamma},\vec{\beta})}$~\cite{fakhimi2023quantum}. Nella pratica questo valore si ottiene eseguendo molto volte il quantum circuit come vedremo più avanti. Oltre al valore dell'energia del sisema la misura dei \emph{qubits} fornisce anche un'assegnazione delle variabili che rappresenta una possibile soluzione del problema.

The QAOA’s performance depends on the number of levels; given optimal parameters of QAOA, the algorithm asymptotically converges to optimal solutions as the number of levels grows\cite{farhi2014quantum}.

\subsection{Ottimizzazione dei parametri}
QAOA requires tuning the parameters to maximize the expected objective value for a given number of levels.
\[
	(\vec{\gamma^*},\vec{\beta^*}) = \argmin_{\vec{\gamma},\vec{\beta}} F(\vec{\gamma},\vec{\beta}) 
\]
This parameter tuning problem is nonconvex and NP-hard~\cite{fakhimi2023quantum}. Even with one level, the gradient of $F(\vec{\gamma},\vec{\beta})$  vanishes rapidly; this phenomenon is known as the barren plateau~\cite{arrasmith2021effect, mcclean2018barren}. The barren plateau expands exponentially as the number of qubits
grows, making the tuning process even more challenging~\cite{mcclean2018barren}.

Nonostante la difficoltà nell'ottimizzare i parametri l'algoritmo QAOA prensenta numerosi aspetti interessanti e la riceca in questo ambito è ancora mlto attiva. Nel nostro lavoro ci da la possibilità di cambiare architettura di rferimento e provare a risolvere un problema in forma QUBO ottenuto con la pipeline di QA-Prolog con quantum gate based computer. Nelle prossime sezioni vedremo come implementare nella pratica l'algoritmo QAOA.

\section{From QAC to QAOA}
Nella sezione preceente ci siamo concetrati sugli aspetti teorici dell'algoritmo QAOA. In questa sezione invece mostriamo come usare il Framework Qiskit di IBM per risolvere un problema di minimizzazione come quelli ottenuti al termine della pipeline di QA-Prolog. Vedremo cioè come implementare praticamente l'algoritmo QAOA.

\subsection{Obbiettivo}
Al termine della pipeline di QA-Prolog otteniamo un problema in forma QUBO, in questa rappresentazione le variabili $x_i$da ottimizzare appatengono all'insieme $\set{0, 1}$. Quando passiamo all'architettura quantum gate based il valore misurato dei \emph{qubits} appartiene all'insieme $\set{-1, 1}$. Dobbiamo quindi passare da una rappesentazione in forma QUBO a un problema in forma Ising. Qiskit inoltre prevede che la matrice $\mathbf{H}$ sia espressa come operatori di Pauli.

Con questi passaggi abbiamo reso compatbile l'output di QA-Prolog con l'input dell'algoritmo QAOA. Per completare l'implementazione dobbiamo poi ottimizzare i parametri definiti nell'Equazione~\ref{} in modo ma minimizzare il valore atteso di $\mathbf{H}$. Infine, dopo aver trovato i parametri ottimi campioneremero il risultato del quantum circuit un'ultima volta per ottenere il risultato del nostro problema.

Per ricapitolare per implementare l'algoritmo QAOA con Qiskit partendo dall'output della pipeline di QA-prolog dobbiamo:
\begin{itemize}
	\item passare da foma QUBO a Forma Ising;
	\item esprimere l'osservabile $\mathbf{H}$ con operatori di Pauli;
	\item ottimizzare i parametri dell'Equazione~\ref{};
	\item campionare il quantum circuit ottimizzato.
\end{itemize}

\subsection{From QUBO to Ising}
Nella Sezione~\ref{} abbiamo visto che una traduzione da problema in forma QUBO a forma Ising esiste sempre. È possibile implementare questa traduzione manualmente seguendo~\cite{bibid}, in questo caso useremo l'implementazione fornita dalla libreria Python \mintinline{text}|qubovert|.

\begin{listing}[h]
	\mylisting{.8}{Python}{code/qaoa/qubo_to_ising.py}
	\caption{text}
	\label{lst:qubo_ising}
\end{listing}

Il Listato~\ref{lst:qubo_ising} mostra come sia possibile passare da una formulazione all'altra tramite la libreria \mintinline{text}|qubovert|:
\begin{enumerate}
	\item costruiamo un problema QUBO avendo come input l'output della pipeline di QA-Prolog (\mintinline{text}|qubo_mat|);
	\item trasformiamo il problema QUBO in forma Ising;
	\item estraiamo il dizionario che assegna i coefficienti numerici $a_i$ nella Formula~\ref{};
	\item scariamo il termine noto;
	\item costriuamo la matrice con i pesi del problema in forma Ising.
\end{enumerate}
L'aspetto più unteressante di questo procedimento è lo scarto del termine noto, questo valore serve per allineare il ground state del problema in Forma QUBO e Ising (Equazione~\ref{}) alla stessa energia. Siccome il nostro obbiettivo trovare il ground state, a prescindere dal suo valore possiamo ignorare questo offset e considerare soli i coefficienti delle variabili.

\subsection{Hamiltonian as Pauli operator}
Consideriamo l'Hamiltoniana $\mathbf{H}$ triangolare superiore che descrive il nostro problema in forma Ising. Ogni elemento appartenente alla diagonale $h_{i,i}$ viene tradotto in un gate che misura la componente $Z$ del \emph{qubit} $i$ e applica la trasformazione identità a tutti gli altri \emph{qubits}\footnote{L'operatore che misura la componente $Z$ è $\sigma_z$ (Equazione~\ref{eqn:pauli_op}), qesto operatore è sia Hermitiano che unitario, quindi descrive sia un'osservabile che un quantum gate valido. Lo stesso vale per l'operatore identità.}. Gli elementi $h_{i,j}$\footnote{Con $j>i$ perchè $\mathbf{H}$ è triangolare superiore.} vengono trasformati in gate che misurano lo spin lungo l'asse $Z$ dei \emph{qubits} $i$ e $j$ e applicano l'operatore identità agli altri \emph{qubits}. 

Seguendo il tutorial di IBM~\cite{ibm_qaoa_tutorial} trasformiamo la matrice che descrive il problema in forma Ising operatori di Pauli. Il Listato~\ref{lst:pauli_op} mostra come sia possibile automatizzare tale riscrittura. Alla linea $4$ generiamo l'osservabile per gli elementi sulla diagonale, alla linea $6$ l'osservabile per gli altri elementi della matrice.
\begin{listing}[h]
	\mylisting{.9}{Python}{code/qaoa/matrix_to_pauli.py}
	\caption{text}
	\label{lst:pauli_op}
\end{listing}

Possiamo infine costruire un osservabile unico a partire dalla lista di operatori di Pauli che abbiamo appena costruito. Per fare questo usaimo una funzione di libreria messa a disposizione da Qiskit. Il Listato~\ref{lst:cost_func} mostra come ottenere la funzione costo, dove \mintinline{text}|pauli_list| sono le gli operatori di Pauli che abbiamo generato con il Listato~\ref{lst:pauli_op} e \mintinline{text}|n_qubits| il numero di \emph{qubits} che descrivono il nustro problema.
\begin{listing}[h]
	\mylistingnonumber{.95}{Python}{code/qaoa/cost_func.py}
	\caption{text}
	\label{lst:cost_func}
\end{listing}

\subsection{Circuito finale}
Come visto nella Sezione~\ref{} l'algoritmo QAOA non utilizza direttamente la funzione costo, ma i gate paraemtrici $555555555$ e $5555555555555555$ ripetuti un cero numero di volte. 

Qiskit mette a disposizione del programmatore un costruttore per generare il quantum circuit che implementi l'algoritmo QAOA; questo costruttore permette di specificare:
\begin{itemize}
	\item funzione costo: quella ottenuta nel Listato~\ref{lst:cost_func};
	\item numero di ripetizioni;
	\item stato iniziale: sotto forma di quantum circuit. Se non viene specificato nessuno stato iniziale si applica a ogni \emph{qubit} la porta di Hadamard per partire da uno stato con soprvrapposizione massima;
	\item operatore di mixing: questo parametro può essere specificato come operatore base oppure come quantum circuit, se non viene specificato nessun parametro si applica il gate Not\ref{}
\end{itemize}

Il circuto termina misurando il valore dei \emph{qubits}. Il Listato~\ref{lst:circ_build} mostra la costruzione del circuito.

\begin{listing}[h]
	\mylisting{.85}{Python}{code/qaoa/circ_build.py}
	\caption{text}
	\label{lst:circ_build}
\end{listing}

\paragraph{Esempio:} Consideriamo un problema descritto da tre \emph{qubits} e due ripetizioni del blocco base di QAOA. La Figura \ref{fig:qaoa_circ} mostra: in~\hyperref[fig:cost_mat]{a} la matrice che descrive il prolema; in \hyperref[fig:qaoa_1]{b} il circuito in versione compatta; in \hyperref[fig:qaoa_2]{c} il circuito espanso per osservare come i coefficienti della matrice sono stati tradotti in osservabili.

\begin{figure}[h]
	\centering
	\begin{subfigure}[c]{0.3\linewidth}
		\centering
		\begin{tikzpicture}
			\node (a) at (0,0) {
				$
				\mathbf{H} = 
				\begin{pmatrix}
					1&0&1\\
					0&1&1\\
					0&0&0
				\end{pmatrix}
				$
			};
		\end{tikzpicture}
		\caption{Text}
		\label{fig:cost_mat}
	\end{subfigure}
	\qquad
	\begin{subfigure}[c]{0.55\linewidth}
		\includegraphics[width=\linewidth]{qaoa/qaoa_1}
		\caption{text}
		\label{fig:qaoa_1}
	\end{subfigure}\\
	\vspace{.5cm}
	\begin{subfigure}[c]{\linewidth}
		\includegraphics[width=\linewidth]{qaoa/qaoa_2}
		\caption{Text}
		\label{fig:qaoa_2}
	\end{subfigure}
	\caption{Text}
	\label{fig:qaoa_circ}
\end{figure}

Quadrando con attenzione gli operatori che costituiscono il gate $5555555555$ (primo e terzo blocco nella Figura~\ref{fig:qaoa_2}) possiamo notare che l'ordine dei qubit è invertito. Ad esempio l'operatore $IIZ$ rappresenta l'elemento di $\mathbf{H}$ in posizione $(0,0)$ e l'operatore $ZZI$ l'elemento in posizione $(1,2)$. Questa considerazione è da tenere presente quando leggeremo i risultati del circuito.

\subsection{Ottimizzazione dei parametri}
Definiamo una funzione per ricavare il costo  $F(\vec{\beta}, \vec{\gamma})$ (Equazione~\ref{}) come valore attteso del risultato dell'esecuzione del quantum circuit definito prima. Questo costo va minimizzato ottimizzando i parametri $\vec{\beta}$ e $ \vec{\gamma}$. Seguendo \cite{ibm_qaoa_tutorial} utilizziamo la libreria Python \mintinline{text}|scipy| per minimizzare $F(\vec{\beta}, \vec{\gamma})$.

\begin{listing}[h]
	\mylisting{.95}{Python}{code/qaoa/cost_func_physic.py}
	\caption{text}
	\label{lst:cost_func_physical}
\end{listing}

Il Listato~\ref{lst:cost_func_physical} mostra come ricavare $F(\vec{\beta}, \vec{\gamma})$ come valore atteso:
\begin{itemize}
	\item alla linea $5$ creiamo un \emph{primitive unified blocks} (\mintinline{text}|pub|): una tripla contenente il circuito da eseguire (\mintinline{text}|ansatz|), l'osservabile da misurare (\mintinline{text}|isa_hamiltonian|) e i paramteri del circuito (\mintinline{text}|params|);
	\item alla linea $6$ stimiamo $F(\vec{\beta}, \vec{\gamma})$ eseguendo multiple run del circuito;
	\item da linea $7$ a $9$ recuperiamo e ritorniamo il risultato.
\end{itemize}
\begin{listing}[h]
	\mylisting{.95}{Python}{code/qaoa/minimize.py}
	\caption{text}
	\label{lst:minimize}
\end{listing}

Il Listato~\ref{lst:minimize} mostra il processo di minimizzazione: creiamo uno stimatore che esegue $1000$ run del circuito (linee $2$ e $3$) e usiamo la funzione di libreria \mintinline{text}|scipy.optimize.minimize|\footnote{La ui documentazione è dispnibile a \url{URL}.} per minimizzare $F(\vec{\beta}, \vec{\gamma})$ ottimizzando i parametri  $\vec{\beta}$ e $ \vec{\gamma}$. Infine creiamo il circuito ottimizzato assegnado al circuito parametrico il valore dei parametri ottimizzati (linea $12$). In questo listato, \mintinline{text}|candidate_circuit| è la versione transpiled\footnote{Sezione~\ref{}} del circuito \mintinline{text}|qc| ottenuto nel Listato~\ref{lst:circ_build}.

\subsection{Recupero dei risultati}
Ora che abbiamo ottimizzato il vlaore dei parametri possiamo, un'ultima volta, fare sampling del quantum circuit, questa volta non per stimare il valore minimo i $F(\vec{\beta}, \vec{\gamma})$, ma per stabilire il valore dei \emph{qubits} che generano tale valore.

\begin{listing}[h]
	\mylisting{.75}{Python}{code/qaoa/result.py}
	\caption{text}
	\label{lst:result}
\end{listing}\

Nel Listato~\ref{lst:result} creiamo un  \mintinline{text}|Sampler| (linea $1$) che run $10000$ volte il circuito ottimizzato (linea $4$). In questo caso usiamo un  \mintinline{text}|Sampler| non un  \mintinline{text}|Estimator| come nel Listato~\ref{lst:minimize} perché siamo interessati al valore dei \emph{qubits} non al valore di un'osservabile. Alla linea $5$ recuperiamo i risultati.

\section{Experiments}
Applichiamo l'algoritmo QAOA per risolvere un problema in forma QUBO ottenuto con la pipeline di QA-Prolog.
\subsection{Problema}
Consideriamo il problema di soddisfacibilità $y = x_1 \land\neg(x_2 \lor x_3)$ (Equazione~\ref{eqn:simple_logic}). Questo problema si presta bene all'applicazione di QAOA perché possiamo leggere la soluzione direttamente dal valore dei \emph{qubits}\footnote{Della possibilità di usare QA-Prolog per rendere leggibile l'output di QAOA parleremo in \ref{}.}. 

Come mostrato nell'Esempio~\ref{} partire dal codice Prolog per risolvere il problema di soddisfacibilità introduce parecchio overhead. Per contenere al massimo il numero di \emph{qubits} necessarei per rappresentare il problema percorriamo solo una parte della pipeline, partendo dal codice Verilog dell'Esempio~\ref{}.

Il problema QUBO che otteniamo è rappresentato da una matrice $11\times11$, cioè abbiamo bisogno di undici \emph{qubits} per rappresentare il nostro problema.

\begin{figure}[h]
	\centering
	\begin{subfigure}[c]{0.437\linewidth}
		\includegraphics[width=\linewidth]{qaoa/h_qubo}
		\caption{Text}
		\label{fig:h_qubo}
	\end{subfigure}
	\hfill
	\begin{subfigure}[c]{0.523\linewidth}
		\includegraphics[width=\linewidth]{qaoa/h_ising}
		\caption{text}
		\label{fig:h_ising}
	\end{subfigure}
	\caption{Text}
	\label{fig:h}
\end{figure}

La Figura~\ref{fig:h} mostra la rappresentazione matriciale del permblema sia in forma Qubo che Ising. Le due rappresentazioni hanno ground state con valore diverso, ma che corrisponde alla stessa configurazione dei \emph{qubits}.

\subsection{Esecuzione}
Per risolvere il probleama di soddisfacibilità costruiamo un QAOA con $10$ livelli. Aumentando il numero di livelli l'algoritmo converge sempre meglio alla soluzione ottima. 

\begin{wrapfigure}{r}{.5\linewidth}
	\centering
	\vspace{-.5cm}
	\includegraphics[width=.95\linewidth]{qaoa/qaoa_opt}
	\caption{text}
	\label{fig:optimization}
\end{wrapfigure}
Aumentare i livelli però comporta anche un tempo di escuzione maggiore e se si esegue su hardware quantistico (e non attraerso simulatori) più è profondo il circuito più si va incontro al fenomeno della decoerenza~\cite{bibid}; inoltre il numero di parametri da ottimizzare è pari al doppio dei livelli, aumentare il numro di livelli quindi comporta anche rendere più complessa l'ottimizzazione dei parametri. La Figura~\ref{fig:optimization} mostra il processo di ottimizzazione dei parametri: man mano che si trovano parametri migliori la stima del costo diminuisce. In questo caso i parametri sono stati ottimzzati dopo $130$ iterazioni circa. Ognuna di queste iperazioni prevede $1000$ run del quantum circuit.

\begin{figure}[h]
	\centering
		\includegraphics[width=.8\linewidth]{qaoa/qaoa_r_1}
	\caption{Text}
	\label{fig:qaoa_r_1}
\end{figure}

Trovati i parametri possimao fare sample del circuito ottimizzato per ottenere la risposta del nostro problema. La Figura~\ref{fig:qaoa_r_1} mostra il risultato del sample ($10000$ run del quantum circuit) evidenziando in viola i quattro risultati più frequenti.

\begin{wrapfigure}{r}{.5\linewidth}
	\centering
	\vspace{-.5cm}
	\includegraphics[width=.95\linewidth]{qaoa/qaoa_r_2}
	\caption{text}
	\label{fig:qaoa_r_2}
\end{wrapfigure}
L'istogramma di Figura~\ref{fig:qaoa_r_1} è di difficile interpretazione, inoltre mostra il riultato di tutti gli undici \emph{qubits} che formano il polinomio da minimizzare. Noi siamo interessati solo al valore delle variabili $x_1$, $x_2$ e $x_3$, le altre sono ancillari. La Figura~\ref{fig:qaoa_r_2} mostra il risutlato del sample in cui consideriamo solo le variabili di nostro interesse. Queste sono state estratte verificando quali \emph{qubits} erano stati mappati da QMASM nelle varibili $x_1$, $x_2$ e $x_3$ e condierando che la bitstring ottenuta è da leggere al contrario

Il risultato ottenuto è quello atteso, infatti abbiamo ottenuto la stringa $1\ 0\ 0$ come sample più probabile, questi valori sono da associare rispettivamente alle variabili $x_1$, $x_3$ e $x_2$ e possiamo verificare facilmente che questo assegnamento soddisfa l'Equazione~\ref{eqn:simple_logic}.

\subsection{Risultati}
Sia il processo di sampling che il processo di ottimizzazione delle variabili introducono dei processi stocastici all'interno dell'esecuzione di QAOA. Verifichiamo quindi se eseguendo più volte l'ottimizzazione dei parametri e il sampling finale otteniamo sempreil risultato atteso\footnote{il codice che automatizza lesecuzione degli esperimenti è disponibie a \url{URL}}.

Siccome il tempo di esecuizione di ciascun esperimento non è trescurabile il numero diesperimenti fatti è contenuto e non permette di fare statistiche attendibili. È però sufficiente per mostrare come l'agoritmo QAOA non sempre sia in grado di trovare la risposta corretta al problema di minimizzazione.

\begin{figure}[h]
	\centering
	\includegraphics[width=.8\linewidth]{qaoa/qaoa_r_3}
	\caption{Text}
	\label{fig:qaoa_r_3}
\end{figure}

La Figura~\ref{fig:qaoa_r_3} mostra i risultati (dei \emph{quits} interessanti) di 12 ripetizioni dell'esperimento. Nell'immagine sono stati evidenziati con lo sfondo verde tutti i casi in cui è stat individuato l'assegnamento corretto delle variabili  $x_1$, $x_2$ e $x_3$, mentre sono evidenziati in rosso i casi in cui il risultato più probabile è sbagliato.

Il umero di esperimenti è molto basso e analisi statistiche avanzate non sarebbero attendibili. Possiamo comunque fare alune osservazioni qualitative. Tipicamente nei casi in cui la soluzione corretta viene trovata la probabilità di quest'ultima è discretamente maggiore di quella degli altri risultati. Quando invece la soluzione è sbagliata i diversi risultati sono distribuiti i maniera abbastanza uniforme.

\section{Conclusion}
In questo capitolo abbiamo presentato un nuovo algoritmo per risolvere un problema di minimizzazione in fomra QUBO. Dal punto di vista teorico abbiamo visto che QAOA ci garantisce un miglioramento monotono delle performance all'aumentare del numero di livelli (a patto di aver trovato dei parametri ottimi), mentre il QAC non garantisce questa monotonia aumentando il tempo di annealing. Abbiamo poi presentato le criticità di QAOA evidenziando come il problema di trovare dei parametri ottimi sia in generale NP-Hard.

La discussione si è poi spsotata sull'implementazione pratica e abbiamo visto come usare il Framework Qiskit per risolvere su macchine quantum gate based di IBM un problema originalmente in forma QUBO. Abbiamo presentato delgi script Python per automatizzare la traduzione, abbiamo visto come recuperare i risultati e come interpretarli. Infine abiamo discusso la qualita delle soluzioni formite da QAOA ripetendo il tunig dei parametri e il sampling del circuto ottimizzato discutendo i risultati ottenuti.

Nel capitolo precedente abbiamo mostrato come aggiungere un passo di riscrttura a monte della pipeline di QA-Prolog. Anche questo capitolo può essere considerato un ulteriore passo di riscrittura di un problema in formati diversi ma equivalenti; questa volta a valle della pipeline abbiamopreso l'output formito da QMASM e lo abbiamo risvritto in un formato compatibile con QAOA.